\section{\heiti 实验与评估}

本节系统评估本文方法的有效性。实验基于真实抽水蓄能电站财务评价模型（数字化系统财务模型边界\_抽水蓄能\_v12.xlsx），依次报告数据集描述、消融实验、知识图谱质量、问答系统评估和案例分析。

\subsection{\heiti 数据集}

实验使用的财务评价模型为某抽水蓄能电站项目的完整财务评价模型，包含14个工作表、57,577个非空单元格、42,109个公式。模型类型为抽水蓄能电站财务评价模型，涵盖投资概算、资金筹措、折旧摊销、成本费用、利润、现金流和资产负债的全链条计算。

\textbf{分层采样策略。} 为评估公式语义分类效果，按各Sheet公式数量比例分层随机采样299条公式（随机种子42），采样比例与各Sheet公式数占总公式数的比例一致。分层采样确保各Sheet的代表性，避免因某些Sheet公式规模大而在测试集中占比过高。

\textbf{金标准构建。} 对采样得到的299条公式，由2名财务建模专家独立标注业务类型（12类），以高置信度（$\geq 0.90$）自动标注结果为基础，人工复核不一致项。最终金标准覆盖12类业务标签，其中融资计算、折旧摊销、成本计算、收入计算和利润计算占比最高，与模型的业务结构一致。

\subsection{\heiti 消融实验}

为验证上下文规则方法的有效性，设计三种分类方法进行消融对比：

\textbf{方法A（纯公式结构）：} 仅使用公式语法特征（函数名、运算符、常量模式）分类，不使用任何上下文信息（Sheet名、行标题、列标题）。

\textbf{方法B（纯上下文规则）：} 仅使用上下文特征（Sheet名、行标题、列标题）进行12类业务规则匹配，不解析公式语法。

\textbf{方法C（上下文规则优先，公式结构补充）：} 上下文规则优先分类，公式结构规则对未覆盖部分进行补充，两者互为fallback。

实验结果如表~\ref{tab:ablation} 所示。

\begin{table}[H]
	\centering
	\caption{语义分类消融实验结果（n=299）}
	\label{tab:ablation}
	\vspace {-2.5mm}
	\begin{center}
		\begin{tabular}{lccc}
			\toprule
			方法 & 准确率/\% & Macro-F1/\% & 覆盖率/\% \\
			\hline
			方法A（纯公式结构） & 20.7 & 12.3 & 37.1 \\
			方法B（纯上下文规则） & 88.3 & 85.6 & 78.3 \\
			方法C（上下文优先+补充） & 88.3 & 85.8 & 78.3 \\
			\bottomrule
		\end{tabular}
	\end{center}
\end{table}

\textbf{发现一：方法A 20.7\%准确率的根因分析。} 方法A（纯公式结构规则）准确率仅20.7\%，覆盖率37.1\%，性能显著低于其他方法。根本原因在于：该模型77.6\%的公式为SUM/IF/算术运算，公式语法本身不携带业务语义——同一SUM公式在不同行可能代表融资计算或折旧摊销，取决于所在行的业务上下文。这从实证角度证明了Excel财务模型公式黑盒问题的本质：语法与语义的结构性解耦，而非单纯的复杂性问题。

\textbf{发现二：方法B/C 88.3\%的机制解释。} 方法B（纯上下文规则）准确率达88.3\%，方法C在此基础上增加公式结构补充后准确率持平。这说明上下文（Sheet名+行标题+列标题）是业务语义的主要载体，公式语法几乎不提供额外的语义区分信息。方法C的公式结构补充主要覆盖了少数含IRR/NPV函数等特殊模式的公式，对整体准确率影响有限。

\textbf{发现三：跨业务Sheet现象。} 剩余11.7\%的错误中，28/35条集中在时间序列Sheet中含折旧关键词的行（如行标题为当年折旧月份数、折旧摊销时间序列等）。追溯根因，该模型的时间序列Sheet内嵌了折旧摊销时间轴的计算逻辑，导致Sheet名称（时间序列）与行标题（含折旧）存在语义冲突——我们称之为跨业务Sheet现象。这类高歧义公式是规则方法的天然边界，也是引入LLM增强分类的优先目标。

\subsection{\heiti 知识图谱质量}

对构建的知识图谱进行完整性验证，统计结果如表~\ref{tab:kg_quality} 所示。

\begin{table}[H]
	\centering
	\caption{知识图谱质量评估}
	\label{tab:kg_quality}
	\vspace {-2.5mm}
	\begin{center}
		\begin{tabular}{lc}
			\toprule
			评估项 & 数值 \\
			\hline
			实体总数 & 42,138 \\
			关系总数 & 65,333 \\
			\quad \textit{contains} 关系 & 14 \\
			\quad \textit{has\_cell} 关系 & 42,109 \\
			\quad \textit{depends\_on} 关系 & 9,146 \\
			\quad \textit{uses\_function} 关系 & 13,547 \\
			\quad \textit{computes} 关系 & 513 \\
			\quad \textit{influences} 关系 & 4 \\
			孤立节点占比 & 17.7\% \\
			图谱连通性 & 弱连通（最大连通分量含99.4\%节点） \\
			\bottomrule
		\end{tabular}
	\end{center}
\end{table}

\subsection{\heiti 问答系统评估}

在13个典型查询模板上测试意图识别准确率，结果如表~\ref{tab:qa_eval} 所示。

\begin{table}[H]
	\centering
	\caption{问答意图识别准确率（n=13）}
	\label{tab:qa_eval}
	\vspace {-2.5mm}
	\begin{center}
		\begin{tabular}{lcc}
			\toprule
			查询类型 & 正确数 & 准确率/\% \\
			\hline
			指标位置查询 & 3/3 & 100.0 \\
			正向追溯 & 3/3 & 100.0 \\
			反向影响 & 3/3 & 100.0 \\
			公式类型查询 & 2/2 & 100.0 \\
			多指标混合查询 & 1/2 & 50.0 \\
			\hline
			综合 & 12/13 & 92.3 \\
			\bottomrule
		\end{tabular}
	\end{center}
	\begin{center}
		\parbox{110mm}{\zihao{5-}\heiti 注：} {\zihao{5-}多指标混合查询中1条误识别，原因为查询文本同时包含"IRR"和"总投资"两个指标词，模板匹配选择优先级冲突。}
	\end{center}
\end{table}

\subsection{\heiti 案例分析：IRR路径追溯}

以IRR（内部收益率）指标为例，演示完整的计算路径追溯。IRR位于表5-利润表-资本金的D105单元格，其正向追溯路径覆盖了从收入、成本到利润的完整计算链：

\begin{enumerate}
	\item IRR(D105) $\leftarrow$ 各期净现金流(D52-D104)
	\item 各期净现金流 $\leftarrow$ 净利润 + 折旧摊销 $-$ 资本支出
	\item 净利润 $\leftarrow$ 营业收入 $-$ 总成本 $-$ 税金
	\item 营业收入 $\leftarrow$ 电价 $\times$ 售电量
	\item 总成本 $\leftarrow$ 运维成本 + 财务费用 + 折旧摊销
\end{enumerate}

通过知识图谱的 \textit{depends\_on} 关系多跳查询，IRR完整计算链可在<0.5秒内返回。相比之下，人工追溯需依次打开14个工作表、定位相关公式、逐个确认依赖关系，耗时约45-60分钟。效率提升超5,000倍。

反向影响分析同样高效：修改参数输入表中的某利率参数，系统沿 \textit{depends\_on} 关系正向追溯，1.2秒内返回该参数影响的1,021个下游单元格，覆盖资金筹措表、折旧摊销表、成本费用表等5个工作表。
